% Generated by build_manuscript.py from 01_introduction.md
\section{1 Introduction}\label{introduction}

Estimating the remaining compressive strength of an impacted composite requires connecting observable damage to a mechanical assessment. Surface appearances, internal damage images and compression-after-impact (CAI) strength measurements describe different aspects of that problem. Public composite-impact datasets make these observations available in corresponding specimen records \citep{hasebe2022data, hasebe2025data}. They also make it possible to ask how much of an internal image is useful for a given assessment task. When internal observations are acquired incrementally, prediction quality depends on both the available evidence and the choice of what to observe next.

A C-scan image can supply spatial information to a strength regressor once that image has been obtained. Recent work in Advanced Engineering Informatics directly predicts CAI strength and impact energy from C-scan images using a convolutional model \citep{mack2026}. This supports a route from image content to mechanical assessment. An additional decision arises when the internal image is not initially available in full: which region should be acquired before the others, and how should that choice change after a measurement? A good complete-input regressor does not by itself specify an acquisition sequence or the quality of predictions along that sequence. Conversely, a spatially economical order is useful for assessment only if the acquired evidence supports the target prediction.

Adaptive inspection addresses parts of this decision problem. Autonomous ultrasonic work has used sequential observations to build damage-indication fields and guide subsequent measurements \citep{fuentes2020}. More generally, dynamic feature-acquisition methods select information according to a prediction objective and an acquisition cost \citep{shim2018, janisch2019, covert2023}. These ideas motivate making the downstream task explicit. For CAI assessment, the objective considered here is the error of the predicted strength over a limited acquisition range, together with its final error. The question is therefore not simply whether an observed region looks damaged, but whether the selected observations form a useful evidence sequence for this particular assessment model.

The available modalities create a natural information hierarchy. A surface image can be available before the internal examination, giving an inexpensive prior within the replay protocol. Internal C-scan content becomes visible only where acquisition has been completed. The next decision can then combine the initial surface information with the current observed subset and strength prediction. This organization requires care: storing complete images for offline evaluation must not allow the policy to use hidden internal content, and using strength labels to train the policy must not expose those labels during acquisition. Without explicit visibility rules, an apparent acquisition advantage could instead reflect information that would not have been available when the decision was made.

We formulate this process as task-driven multimodal C-scan acquisition. A frozen surface vision-language model (VLM) proposes candidate regions and ordinal confidence, which provide an initial spatial prior. A learned actor makes the acquisition decisions, updating its state after each newly observed internal cell. A separate frozen CAI predictor converts the currently available surface and internal evidence into a strength estimate. The VLM is therefore a source of surface priors rather than the controller of each step. The learning objective combines normalized error area across the acquisition budget with terminal absolute error, so it distinguishes early useful observations from the same prediction change obtained later.

To make the resulting comparison interpretable, all evaluated acquisition policies share the same downstream predictor and legal whole-cell budget rule. Fixed orders, a learned shared ranking, a specimen-conditioned open-loop policy and feedback policies represent different ways to organize observation. Additional comparisons remove the VLM prior or replace spatial interactions with mean feedback. Complete internal input provides a separate quality reference at full acquisition. This arrangement asks where state dependence helps in the observed trajectories while retaining cases in which an added component does not improve the estimate. It avoids attributing all differences between separately trained policies to a single architectural factor.

The study makes three contributions. First, it specifies a sequential CAI-assessment task with explicit surface availability, hidden internal regions, native-pixel acquisition costs and a common prediction model. Second, it develops an acquisition policy that conditions successive choices on permitted multimodal evidence and is trained by the downstream trajectory loss rather than an action demonstration. Third, it characterizes quality under matched acquisition caps, acquisition requirements at matched empirical quality and the timing-weighted error changes in the saved trajectories. The timing decomposition is an algebraic explanation of the loss, not a causal test of damage mechanisms.

We evaluate this framework by offline replay on 50 validation specimens from six source domains, with checkpoints selected on that same validation set. The paper therefore reports an exploratory within-cohort comparison, not independent deployment validation. Results address the linked questions of same-cap prediction quality, shared versus adaptive ordering, the stages at which error changes contribute, and the remaining gap to complete input. This evidence supports a concrete assessment of the acquisition method while keeping physical scan time, autonomous stopping and generalization to new specimen populations as distinct questions beyond the evaluated protocol.
